\documentclass[10pt,a4paper]{amsart}
\usepackage{amsmath,amssymb,amsthm,graphicx}
\usepackage{verbatim}
\newenvironment{reference}{\comment}{\endcomment}
%\newenvironment{reference}{}{}
\newenvironment{slogan}{\comment}{\endcomment}
\newenvironment{history}{\comment}{\endcomment}
% For commutative diagrams we use Xy-pic
\usepackage[all]{xy}

% We use 2cell for 2-commutative diagrams.
\xyoption{2cell}
% Theorem environments.
%
\theoremstyle{plain}
\newtheorem{theorem}[subsection]{Theorem}
\newtheorem{proposition}[subsection]{Proposition}
\newtheorem{lemma}[subsection]{Lemma}

\theoremstyle{definition}
\newtheorem{definition}[subsection]{Definition}
\newtheorem{example}[subsection]{Example}
\newtheorem{exercise}[subsection]{Exercise}
\newtheorem{situation}[subsection]{Situation}

\theoremstyle{remark}
\newtheorem{remark}[subsection]{Remark}
\newtheorem{remarks}[subsection]{Remarks}

\numberwithin{equation}{subsection}

% Macros
%
\def\lim{\mathop{\mathrm{lim}}\nolimits}
\def\colim{\mathop{\mathrm{colim}}\nolimits}
\def\Spec{\mathop{\mathrm{Spec}}}
\def\Hom{\mathop{\mathrm{Hom}}\nolimits}
\def\Ext{\mathop{\mathrm{Ext}}\nolimits}
\def\SheafHom{\mathop{\mathcal{H}\!\mathit{om}}\nolimits}
\def\SheafExt{\mathop{\mathcal{E}\!\mathit{xt}}\nolimits}
\def\Sch{\mathit{Sch}}
\def\Mor{\mathop{\mathrm{Mor}}\nolimits}
\def\Ob{\mathop{\mathrm{Ob}}\nolimits}
\def\Sh{\mathop{\mathit{Sh}}\nolimits}
\def\NL{\mathop{N\!L}\nolimits}
\def\CH{\mathop{\mathrm{CH}}\nolimits}
\def\proetale{{pro\text{-}\acute{e}tale}}
\def\etale{{\acute{e}tale}}
\def\QCoh{\mathit{QCoh}}
\def\Ker{\mathop{\mathrm{Ker}}}
\def\Im{\mathop{\mathrm{Im}}}
\def\Coker{\mathop{\mathrm{Coker}}}
\def\Coim{\mathop{\mathrm{Coim}}}

% Boxtimes
%
\DeclareMathSymbol{\boxtimes}{\mathbin}{AMSa}{"02}

%
% Macros for moduli stacks/spaces
%
\def\QCohstack{\mathcal{QC}\!\mathit{oh}}
\def\Cohstack{\mathcal{C}\!\mathit{oh}}
\def\Spacesstack{\mathcal{S}\!\mathit{paces}}
\def\Quotfunctor{\mathrm{Quot}}
\def\Hilbfunctor{\mathrm{Hilb}}
\def\Curvesstack{\mathcal{C}\!\mathit{urves}}
\def\Polarizedstack{\mathcal{P}\!\mathit{olarized}}
\def\Complexesstack{\mathcal{C}\!\mathit{omplexes}}
% \Pic is the operator that assigns to X its picard group, usage \Pic(X)
% \Picardstack_{X/B} denotes the Picard stack of X over B
% \Picardfunctor_{X/B} denotes the Picard functor of X over B
\def\Pic{\mathop{\mathrm{Pic}}\nolimits}
\def\Picardstack{\mathcal{P}\!\mathit{ic}}
\def\Picardfunctor{\mathrm{Pic}}
\def\Deformationcategory{\mathcal{D}\!\mathit{ef}}
\usepackage{cleveref}
\crefname{originallemma}{original lemma}{original lemmas}
\Crefname{originallemma}{Original lemma}{Original lemmas}
\makeatletter
\newlabel{lemma-injective-dimension-over-polynomial-ring}{{lemma-injective-dimension-over-polynomial-ring}{source}{}{sourcepointer-1}{}}
\newlabel{lemma-injective-dimension-over-polynomial-ring@cref}{{[originallemma][0][]lemma-injective-dimension-over-polynomial-ring}{[1][0][]source}{}{}{}}
\makeatother
\crefname{originalreference}{original reference}{original references}
\Crefname{originalreference}{Original reference}{Original references}
\makeatletter
\newlabel{more-algebra-lemma-base-change-RHom}{{more-algebra-lemma-base-change-RHom}{source}{}{sourcepointer-2}{}}
\newlabel{more-algebra-lemma-base-change-RHom@cref}{{[originalreference][0][]more-algebra-lemma-base-change-RHom}{[1][0][]source}{}{}{}}
\makeatother
\crefname{originalreference}{original reference}{original references}
\Crefname{originalreference}{Original reference}{Original references}
\makeatletter
\newlabel{more-algebra-lemma-finite-injective-dimension}{{more-algebra-lemma-finite-injective-dimension}{source}{}{sourcepointer-3}{}}
\newlabel{more-algebra-lemma-finite-injective-dimension@cref}{{[originalreference][0][]more-algebra-lemma-finite-injective-dimension}{[1][0][]source}{}{}{}}
\makeatother
\title{Extending a dualizing complex along a polynomial ring stays dualizing}
\author{The Stacks Project Authors}
\date{stacks-project:dualizing.tex:local-snapshot-20261012; GNU FDL 1.2 or later}
\begin{document}
\maketitle
\noindent\textit{Editorial note.} All delivered mathematical text is the author's own, copied from the pinned Stacks Project chapter source sources/stacks-dualizing/pinned-original-source.tex; no statement or proof step was rewritten, reordered or removed.\par\medskip
\noindent\textit{Editorial note.} Preamble provenance: the Stacks Project keeps one single shared preamble for every chapter; this card takes it directly from the staged sources/stacks-dualizing/preamble.tex as three verbatim spans (lines 8-12, 14-18 and 154-217), so no derived file was written and nothing was added or redefined. The wrapper class and the base packages are supplied by the assembler.\par\medskip
\noindent\textit{Editorial note.} Bibliography provenance: neither the counted statement nor its proof contains a citation (this lane rescans the whole span before writing the card), so the card carries no bibliography and no entry was reconstructed.\par\medskip
\noindent\textit{Editorial note.} Selection provenance: this card was opened under the file-based control contract (CONTROL.md), with the mechanical pre-card gate run on its statement and proof interval before the card was written.\par\medskip
\noindent\textit{Allowed standard prerequisites (editorial).}\begin{itemize}
\item The finiteness of the injective dimension of the extension of scalars of the dualizing complex. Invoked at the start of the counted proof (source lines 3075-3077). Establishes the first defining property of a dualizing complex for the polynomial ring.
\item The base change behaviour of the derivation of the internal hom along a flat map. Invoked at the end of the counted proof (source lines 3084-3086). Turns the duality quasi-isomorphism over the source into the corresponding quasi-isomorphism over the polynomial ring.
\end{itemize}\medskip
\ifdefined\hypertarget\hypertarget{sourcepointer-1}{}\fi
\noindent\textit{Source reference.} Original lemma lemma-injective-dimension-over-polynomial-ring: sources/stacks-dualizing/pinned-original-source.tex, source lines 591--591. Navigation pointer only: the counted text refers to the Stacks item lemma-injective-dimension-over-polynomial-ring, which is an existing labeled result of this same chapter and is not part of this card. The pointer carries the target's own label as its number and its definition line as its span; its content is not reproduced here.\par\medskip
\ifdefined\hypertarget\hypertarget{sourcepointer-2}{}\fi
\noindent\textit{Source reference.} Original reference more-algebra-lemma-base-change-RHom: sources/stacks-dualizing/pinned-original-source.tex, source lines 2719--2719. Navigation pointer only: the counted text refers to the Stacks item more-algebra-lemma-base-change-RHom, which belongs to another chapter of the project and is not part of this card. The pointer carries the target's own label as its number because this chapter's staged snapshot does not contain the other chapter's tag table; its content is not reproduced here.\par\medskip
\ifdefined\hypertarget\hypertarget{sourcepointer-3}{}\fi
\noindent\textit{Source reference.} Original reference more-algebra-lemma-finite-injective-dimension: sources/stacks-dualizing/pinned-original-source.tex, source lines 3076--3076. Navigation pointer only: the counted text refers to the Stacks item more-algebra-lemma-finite-injective-dimension, which belongs to another chapter of the project and is not part of this card. The pointer carries the target's own label as its number because this chapter's staged snapshot does not contain the other chapter's tag table; its content is not reproduced here.\par\medskip
\begin{lemma}
\label{lemma-dualizing-polynomial-ring}
Let $A$ be a Noetherian ring. If $\omega_A^\bullet$ is a dualizing
complex, then $\omega_A^\bullet \otimes_A A[x]$ is a dualizing
complex for $A[x]$.
\end{lemma}

\begin{proof}
Set $B = A[x]$ and $\omega_B^\bullet = \omega_A^\bullet \otimes_A B$.
It follows from Lemma \ref{lemma-injective-dimension-over-polynomial-ring}
and More on Algebra, Lemma \ref{more-algebra-lemma-finite-injective-dimension}
that $\omega_B^\bullet$ has finite injective dimension.
Since $H^i(\omega_B^\bullet) = H^i(\omega_A^\bullet) \otimes_A B$
by flatness of $A \to B$ we see that $\omega_A^\bullet \otimes_A B$
has finite cohomology modules. Finally, the map
$$
B \longrightarrow R\Hom_B(\omega_B^\bullet, \omega_B^\bullet)
$$
is a quasi-isomorphism as formation of internal hom commutes with
flat base change in this case, see
More on Algebra, Lemma \ref{more-algebra-lemma-base-change-RHom}.
\end{proof}

\end{document}
